\par\Needspace{7\baselineskip}\bigskip
\begin{center}
\pdfbookmark[2]{Gauss an Olbers, Göttingen, 31. December 1814}{pallas-brief-18141231}
\textsc{Gauss} an \textsc{Olbers}, Göttingen, 31. December 1814.
\end{center}
\noindent\hspace*{\parindent}. . . . .\quad Eine kurze Nachricht über die letzte Pallasopposition haben Sie wahrscheinlich in unsern Gel. Anz. gelesen. Ich habe seitdem die Elemente noch schärfer angepasst und finde folgende Unterschiede:
\begin{center}
\begingroup\small\setlength{\tabcolsep}{5pt}
\begin{tabular}{@{}lrrrr@{}}
&\multicolumn{2}{c}{Alte Jupitersmasse \(\left[\frac{1}{1067{,}09}\right]\)}&\multicolumn{2}{c}{Jupitersmasse um \(\frac{1}{64}\) vermehrt}\\[3pt]
&\multicolumn{1}{c}{Mittl. Länge}&\multicolumn{1}{c}{Hel. Br.}&\multicolumn{1}{c}{Mittl. Länge}&\multicolumn{1}{c}{Hel. Br.}\\
I.&\(-2''{,}6\)&\(-5''{,}1\)&\(+4''{,}4\)&\(-3''{,}0\)\\
II.&\(-11{,}2\)&\(-12{,}2\)&\(+1{,}8\)&\(-8{,}5\)\\
III.&\(-17{,}3\)&\(-0{,}5\)&\(-5{,}9\)&\(-3{,}3\)\\
IV.&\(+3{,}2\)&\(-6{,}7\)&\(-10{,}8\)&\(-6{,}1\)\\
V.&\(+22{,}8\)&\(+9{,}0\)&\(+3{,}3\)&\(+11{,}5\)\\
VI.&\(+23{,}8\)&\(-3{,}4\)&\(+1{,}1\)&\(-0{,}5\)\\
VII.&\(+16{,}8\)&\(-7{,}2\)&\(+13{,}2\)&\(-5{,}4\)\\
VIII.&\(-17{,}9\)&\(+4{,}5\)&\(-1{,}2\)&\(+2{,}2\)\\
IX.&\(-9{,}6\)&\(-0{,}1\)&\(-4{,}0\)&\(-0{,}7\)\\
X.&\(-7{,}7\)&\(+2{,}2\)&\(-2{,}0\)&\(+5{,}5\)\\
Summa Quadr.&2284&391&381&325.
\end{tabular}
\endgroup
\end{center}
\par\SourceLeaf{432}Hiebei liegen aber nur die Jupitersstörungen nach derjenigen Rechnung, die ich die \textls[180]{specielle} nenne (von einem Jahre zum andern durch mechanische Quadratur fortgesetzt) [zum Grunde]. Die Berechnung nach meinen allgemeinen Formeln, sowie die der Saturnsstörungen, ist noch nicht vollendet. Darüber gehen jedesmal einige Monate hin. Ich werde nun bald meine Methode ausarbeiten und auf eine oder andere Art bekannt machen, um endlich dieser zwar durch guten Erfolg belohnten, aber an sich gar zu lästigen Arbeit quitt zu werden.
\par
